% =============================================================================
% 03-system-calls.tex — Week 3:
% System Calls: Mechanics, Parameter Passing, and Types
%
% Course: Operating Systems (PCC CS-502)
% Anchor reading: Silberschatz, Galvin & Gagne, Operating System Concepts,
%                 8th ed. — Ch. 2 §2.3 (system calls, pp. 55-58) and §2.4
%                 (types of system calls, pp. 58-66); page-verified via
%                 master_supporting_docs/CS502/supporting_books/SilberschatzOSC/index.md.
%                 Bovet & Cesati, Understanding the Linux Kernel, 3rd ed. —
%                 Ch. 10 (system calls, pp. 398-419); page-verified via
%                 .../LinuxKernel/index.md.
%
% Compile with: /compile-latex CS502/03-system-calls
% =============================================================================

\documentclass[aspectratio=169]{beamer}
\input{header}
\coursecode{CS502}

\title{System Calls}
\subtitle{Week 3 --- Mechanics, Parameter Passing, and Types}
\author[Dr. Auqib Hamid Lone]{%
  \textbf{Dr. Auqib Hamid Lone}\\
  \small Assistant Professor in Computer Science \& Engineering
}
\institute[GCET Ganderbal]{%
  \small\textbf{PCC CS-502: Operating Systems}\\[0.3em]
  Department of Computer Science \& Engineering\\
  Government College of Engineering \& Technology (GCET), Ganderbal
}
\date[\today]{\small\today}

\begin{document}

\frame{\titlepage}

% =============================================================================
% ACT 0 --- THE COMMAND THAT KEEPS KNOCKING
% =============================================================================
\section{The Command That Keeps Knocking}

\begin{frame}{Welcome: How Many Times Does \texttt{cp} Enter the Kernel?}
  \begin{block}{The big question}
    You run \texttt{cp notes.txt backup.txt}. The command opens one file,
    creates another, copies every byte, and closes both. How many times does
    the program cross from user mode into the kernel --- and what does it say
    each time?
  \end{block}
  \begin{itemize}
    \item \texttt{cp} is ordinary user code. It cannot touch the disk, the
      page tables, or another process's memory.
    \item Every privileged operation must be requested through a
      \key{system call}.
    \item The request has a name, a number, arguments, and a return value.
  \end{itemize}
  \vspace{0.15cm}
  \begin{alertblock}{Running thread for today}
    We follow \texttt{cp} through its system calls and open up the mechanism
    that carries each request across the boundary.
  \end{alertblock}
\end{frame}

\begin{frame}{Roadmap: Four Questions About the Doorway}
  \begin{enumerate}
    \item \key{What is a system call?} --- the kernel's programming interface.
    \item \key{How does control cross the boundary?} --- the trap and the
      system-call table.
    \item \key{How do arguments travel?} --- registers, memory blocks, and the
      stack.
    \item \key{What kinds of calls exist?} --- the six standard categories.
  \end{enumerate}
  \vspace{0.2cm}
  \muted{The mechanism is hardware; the interface is software; the taxonomy is
  the vocabulary you will use for the rest of the course.}
\end{frame}

\transitionslide{Programmers rarely invoke a system call directly. They call a library function --- and the library eventually executes the trap.}

% =============================================================================
% ACT 1 --- API VS. SYSTEM CALL
% =============================================================================
\section{API vs.\ System Call}

\begin{frame}{Two Views of the Same Service}
  \centering
  \renewcommand{\arraystretch}{1.35}
  \begin{tabular}{p{0.44\textwidth}p{0.44\textwidth}}
    \textbf{Programmer's view (API)} & \textbf{Kernel's view (system call)} \\
    \hline
    \texttt{fopen}, \texttt{fread}, \texttt{fwrite}, \texttt{fclose} &
    \texttt{open}, \texttt{read}, \texttt{write}, \texttt{close} \\
    Standard C library; buffered; portable &
    Kernel entry points; unbuffered; architecture-specific \\
    Returns a \texttt{FILE *} or a count &
    Returns a file descriptor or $-1$ with \texttt{errno} \\
  \end{tabular}
  \vspace{0.35cm}
  \begin{itemize}
    \item The API is a \emph{convenience layer}; the system call is the
      \emph{actual request} to the kernel.
    \item One library call may issue zero, one, or many system calls
      (buffering changes the count).
  \end{itemize}
  \muted{Silberschatz et al., Ch.~2 \S2.3, pp.~55--58.}
\end{frame}

\begin{frame}{What Exactly Is a System Call?}
  \begin{block}{Definition}
    A \key{system call} is the programmatic request from a user program to a
    kernel service. It is the boundary of the operating system's programming
    interface.
  \end{block}
  \begin{itemize}
    \item It looks like a function call in C, but it is \emph{not} an ordinary
      function call.
    \item It changes the privilege level: the CPU enters kernel mode, runs
      kernel code, and returns to user mode.
    \item The kernel executes it on behalf of the calling process, with the
      kernel's privileges --- after validating the request.
  \end{itemize}
  \vspace{0.15cm}
  \muted{The same word is used for the API function, the trap instruction, and
  the kernel routine; context tells you which layer is meant.}
\end{frame}

\begin{frame}{Why Not Link the Kernel Into Every Program?}
  \begin{itemize}
    \item \textbf{Protection:} user code would run with kernel privileges and
      could corrupt the machine.
    \item \textbf{Isolation:} one buggy program could overwrite another's
      memory or the kernel's data structures.
    \item \textbf{Portability:} the same API can sit on different kernels and
      instruction sets.
    \item \textbf{Evolution:} the kernel can change its internals without
      recompiling every application.
  \end{itemize}
  \vspace{0.2cm}
  \begin{alertblock}{The trade-off}
    Every system call costs a mode switch. The design goal is to keep the
    boundary narrow and the calls efficient.
  \end{alertblock}
\end{frame}

\begin{frame}{Socratic Check 1}
  \begin{block}{Question}
    \texttt{printf} is not a system call. What does it eventually call, and
    why does the C library buffer the output first?
  \end{block}
  \vspace{0.25cm}
  \begin{itemize}
    \item \texttt{printf} formats into a user-space buffer.
    \item When the buffer fills or is flushed, the library calls
      \texttt{write}, which traps into the kernel.
    \item Buffering reduces the number of expensive mode switches.
  \end{itemize}
\end{frame}

% =============================================================================
% ACT 2 --- MECHANICS
% =============================================================================
\section{Mechanics of a System Call}

\transitionslide{The trap instruction is the hinge: one instruction moves the CPU from user mode to kernel mode and hands it a number.}

\begin{frame}{The Trap Path}
  \centering
  \begin{tikzpicture}[every node/.style={font=\small}]
    \node[dag-node, minimum width=2.8cm] (user) at (-5.2,0) {User code};
    \node[dag-node, minimum width=2.8cm] (lib) at (-1.8,0) {Library stub};
    \node[dag-node, minimum width=2.8cm] (trap) at (1.8,0) {Trap};
    \node[dag-node, minimum width=2.8cm] (table) at (5.2,0) {Syscall table};
    \node[dag-node, minimum width=2.8cm] (svc) at (5.2,-2.0) {Kernel handler};
    \draw[observed-edge] (user) -- (lib);
    \draw[observed-edge] (lib) -- (trap);
    \draw[observed-edge] (trap) -- (table);
    \draw[observed-edge] (table) -- (svc);
    \draw[observed-edge] (svc) -- ++(-5.2,0) -| (user);
  \end{tikzpicture}
  \vspace{0.2cm}
  \begin{itemize}
    \item The stub places the \key{system-call number} and arguments in
      registers.
    \item The trap saves the user PC and flags, switches to kernel mode, and
      jumps to the kernel's trap entry point.
    \item The kernel indexes the \key{system-call table} by number and calls
      the handler.
  \end{itemize}
\end{frame}

\begin{frame}{Worked Trace: \texttt{write(fd, buf, n)}}
  \begin{enumerate}
    \item \texttt{fd}, \texttt{buf}, and \texttt{n} are placed in the
      argument registers.
    \item The system-call number for \texttt{write} is placed in a designated
      register.
    \item The stub executes \texttt{syscall} (x86-64) or \texttt{int 0x80}
      (legacy x86).
    \item Hardware switches to kernel mode, saves the return address, and
      enters the trap handler.
    \item The kernel validates \texttt{fd}, checks that \texttt{buf} lies in
      user memory, and performs the write.
    \item The return value is placed in a register; the kernel restores user
      mode and returns.
    \item On error, the library sets \texttt{errno} and returns $-1$.
  \end{enumerate}
  \vspace{0.1cm}
  \muted{Bovet \& Cesati, Ch.~10, pp.~398--419; Silberschatz et al., Ch.~2
  \S2.3.}
\end{frame}

\begin{frame}{The System-Call Table}
  \begin{itemize}
    \item A kernel array indexed by system-call number.
    \item Each entry is a function pointer to the handler.
    \item The number is the \key{ABI}: it must stay stable across kernel
      versions.
    \item Adding a call means adding a table entry and a handler --- not
      changing the trap instruction.
  \end{itemize}
  \vspace{0.2cm}
  \begin{block}{Why a table?}
    The trap instruction carries only a number. The table turns that number
    into the right kernel routine without a long search.
  \end{block}
  \muted{On Linux x86-64, \texttt{write} is system call 1; the table lives in
  the kernel's trap entry code.}
\end{frame}

\begin{frame}{Return Values and \texttt{errno}}
  \begin{itemize}
    \item The kernel returns a value in a register: a count, a file
      descriptor, or an error code.
    \item A negative return value signals an error; the library translates it
      into \texttt{-1} and sets \texttt{errno}.
    \item The caller must check the return value --- the kernel cannot print a
      message for it.
    \item The kernel never trusts user pointers or lengths; it validates them
      before copying.
  \end{itemize}
  \vspace{0.2cm}
  \begin{exampleblock}{Example}
    \texttt{read} returns the number of bytes read, \texttt{0} at end of file,
    or $-1$ on error with \texttt{errno} set to \texttt{EFAULT},
    \texttt{EBADF}, or similar.
  \end{exampleblock}
\end{frame}

% =============================================================================
% ACT 3 --- PARAMETER PASSING
% =============================================================================
\section{Parameter Passing}

\transitionslide{A system call is a function call across a privilege boundary. The arguments must cross too --- and the kernel must not trust any of them.}

\begin{frame}{Three Ways to Pass Arguments}
  \centering
  \begin{tikzpicture}[every node/.style={font=\small}]
    \node[dag-node, minimum width=3.4cm, minimum height=1.1cm] (reg) at (-5.0,0)
      {Registers};
    \node[dag-node, minimum width=3.4cm, minimum height=1.1cm] (block) at (0,0)
      {Block / table\\in memory};
    \node[dag-node, minimum width=3.4cm, minimum height=1.1cm] (stack) at (5.0,0)
      {Stack};
    \node[text=primary-gold, align=center] at (-5.0,-1.25) {fast\\few arguments};
    \node[text=primary-gold, align=center] at (0,-1.25) {many arguments\\one pointer};
    \node[text=primary-gold, align=center] at (5.0,-1.25) {natural for\\C calling convention};
  \end{tikzpicture}
  \vspace{0.35cm}
  \begin{itemize}
    \item Silberschatz describes all three; real ABIs choose one or a mix.
    \item The kernel must know where to find the arguments before it can
      validate them.
  \end{itemize}
\end{frame}

\begin{frame}{Method 1: Registers}
  \begin{itemize}
    \item The caller places the first few arguments in designated registers.
    \item Fast: no memory access beyond the registers already saved on entry.
    \item Limited by the number of argument registers.
  \end{itemize}
  \vspace{0.2cm}
  \begin{exampleblock}{Linux x86-64 \texttt{read(fd, buf, count)}}
    \begin{tabular}{@{}ll@{}}
      \texttt{rax} & system-call number for \texttt{read} \\
      \texttt{rdi} & \texttt{fd} \\
      \texttt{rsi} & \texttt{buf} (user virtual address) \\
      \texttt{rdx} & \texttt{count} \\
    \end{tabular}
  \end{exampleblock}
  \muted{The kernel still checks that \texttt{buf} is a legal user address.}
\end{frame}

\begin{frame}{Method 2: Block or Table in Memory}
  \begin{itemize}
    \item The caller stores the arguments in a memory block.
    \item A register holds the \emph{address} of the block.
    \item The kernel copies the block into kernel space and validates every
      field.
  \end{itemize}
  \vspace{0.2cm}
  \begin{exampleblock}{A call with many arguments}
    \texttt{open(path, flags, mode)} can be passed as a table holding the
    pointer to \texttt{path}, \texttt{flags}, and \texttt{mode}; the kernel
    reads the table once.
  \end{exampleblock}
  \muted{The pointer itself is untrusted: the kernel must verify that the whole
  block lies in the caller's address space.}
\end{frame}

\begin{frame}{Method 3: Stack}
  \begin{itemize}
    \item The caller pushes arguments onto the user stack.
    \item The kernel reads them from the saved user stack.
    \item Natural for languages and calling conventions that already pass
      arguments on the stack.
  \end{itemize}
  \vspace{0.2cm}
  \begin{block}{The common requirement}
    Whichever method is used, the kernel must (1) locate the arguments,
    (2) copy them into kernel memory, and (3) validate every pointer and
    length before using it.
  \end{block}
\end{frame}

\begin{frame}{Pointer Validation: Never Trust User Memory}
  \begin{itemize}
    \item A user pointer may point to kernel memory, another process's memory,
      or unmapped memory.
    \item The kernel checks the range against the process's address space.
    \item It copies data with special routines
      (\texttt{copy\_from\_user}/\texttt{copy\_to\_user} on Linux) that handle
      faults safely.
    \item A bad pointer returns \texttt{EFAULT}; it does not crash the kernel.
  \end{itemize}
  \vspace{0.2cm}
  \begin{alertblock}{The security rule}
    Every argument is untrusted input. Validation happens in kernel mode,
    before the kernel acts on the request.
  \end{alertblock}
\end{frame}

\begin{frame}{Socratic Check 2}
  \begin{block}{Question}
    A user program passes a pointer to a kernel address to \texttt{read}. What
    should happen, and why can the kernel not simply dereference it?
  \end{block}
  \vspace{0.25cm}
  \begin{itemize}
    \item Dereferencing it would let user code read or overwrite kernel data.
    \item The kernel checks that the range lies in the user address space.
    \item If it does not, the call fails with \texttt{EFAULT}.
  \end{itemize}
\end{frame}

% =============================================================================
% ACT 4 --- TYPES OF SYSTEM CALLS
% =============================================================================
\section{Types of System Calls}

\transitionslide{System calls are not a random list. They fall into six categories, and every command you run is a combination of them.}

\begin{frame}{The Six Categories}
  \centering
  \footnotesize
  \renewcommand{\arraystretch}{1.3}
  \begin{tabular}{p{0.24\textwidth}p{0.38\textwidth}p{0.30\textwidth}}
    \textbf{Category} & \textbf{Purpose} & \textbf{Typical calls} \\
    \hline
    Process control & create, terminate, wait, load & \texttt{fork}, \texttt{exec},
    \texttt{wait}, \texttt{exit} \\
    File management & create, open, read, write, close & \texttt{open},
    \texttt{read}, \texttt{write}, \texttt{close}, \texttt{lseek} \\
    Device management & request/release, read/write, control & \texttt{ioctl},
    \texttt{read}, \texttt{write} \\
    Information maintenance & get/set time, process, file attributes &
    \texttt{getpid}, \texttt{alarm}, \texttt{sleep}, \texttt{stat} \\
    Communication & create/delete connections, send/receive & \texttt{pipe},
    \texttt{mmap}, \texttt{socket}, \texttt{send}, \texttt{recv} \\
    Protection & set/get permissions & \texttt{chmod}, \texttt{umask},
    \texttt{chown}, \texttt{setuid} \\
  \end{tabular}
  \vspace{0.2cm}
  \muted{Silberschatz et al., Ch.~2 \S2.4, pp.~58--66.}
\end{frame}

\begin{frame}{Process Control}
  \begin{itemize}
    \item \texttt{fork}: create a child process that is a copy of the parent.
    \item \texttt{exec}: replace the current process image with a new program.
    \item \texttt{wait}: block until a child terminates and collect its status.
    \item \texttt{exit}: terminate the calling process and release its
      resources.
  \end{itemize}
  \vspace{0.2cm}
  \begin{exampleblock}{The shell's job}
    A shell \texttt{fork}s, the child \texttt{exec}s the command, and the
    parent \texttt{wait}s. \texttt{cp} itself is the child after \texttt{exec}.
  \end{exampleblock}
\end{frame}

\begin{frame}{File Management}
  \begin{itemize}
    \item \texttt{open}: create or open a file, returning a file descriptor.
    \item \texttt{read}/\texttt{write}: transfer bytes between user memory and
      the file.
    \item \texttt{lseek}: move the file offset without transferring data.
    \item \texttt{close}: release the descriptor and flush state.
    \item \texttt{stat}: read metadata (size, owner, permissions, times).
  \end{itemize}
  \vspace{0.2cm}
  \muted{The file descriptor is a small integer index into the process's
  per-process open-file table --- not the file itself.}
\end{frame}

\begin{frame}{Device, Information, Communication, Protection}
  \begin{columns}[T]
    \begin{column}{0.48\textwidth}
      \begin{block}{Device management}
        \texttt{ioctl} controls device-specific behavior; \texttt{read} and
        \texttt{write} move data. Devices are often accessed through the same
        descriptor interface as files.
      \end{block}
      \begin{block}{Information maintenance}
        \texttt{getpid}, \texttt{alarm}, \texttt{sleep}, \texttt{time};
        \texttt{stat} for file metadata.
      \end{block}
    \end{column}
    \begin{column}{0.48\textwidth}
      \begin{block}{Communication}
        \texttt{pipe}, \texttt{mmap}, \texttt{socket}, \texttt{send},
        \texttt{recv}: move data between processes or across a network.
      \end{block}
      \begin{block}{Protection}
        \texttt{chmod}, \texttt{umask}, \texttt{chown}, \texttt{setuid}: set or
        query access rights.
      \end{block}
    \end{column}
  \end{columns}
\end{frame}

\begin{frame}{Worked Trace: \texttt{cp} Across the Categories}
  \begin{enumerate}
    \item \textbf{Process control:} the shell \texttt{fork}s and the child
      \texttt{exec}s \texttt{cp}.
    \item \textbf{File management:} \texttt{open("notes.txt", O\_RDONLY)} and
      \texttt{open("backup.txt", O\_WRONLY|O\_CREAT)}.
    \item \textbf{File management:} a loop of \texttt{read} into a buffer and
      \texttt{write} to the destination.
    \item \textbf{File management:} \texttt{close} both descriptors.
    \item \textbf{Process control:} \texttt{exit(0)}; the shell \texttt{wait}s.
  \end{enumerate}
  \vspace{0.15cm}
  \muted{One shell command, ten or more system calls. The categories are not
  abstract: they are the steps in the trace.}
\end{frame}

\begin{frame}{Socratic Check 3}
  \begin{block}{Question}
    Which category does \texttt{fork} belong to, and why is it not file
    management even though the child may open files?
  \end{block}
  \vspace{0.25cm}
  \begin{itemize}
    \item \texttt{fork} creates a new process --- process control.
    \item The child inherits the parent's open-file table, but \texttt{fork}
      itself does not create or access a file.
    \item The category describes the \emph{operation}, not the resources the
      process may use afterward.
  \end{itemize}
\end{frame}

% =============================================================================
% ACT 5 --- SYSTEM CALLS IN PRACTICE
% =============================================================================
\section{System Calls in Practice}

\transitionslide{The interface is stable, but every call has a cost. Good systems programmers batch calls and know when the boundary is being crossed.}

\begin{frame}{Numbers, ABI, and Compatibility}
  \begin{itemize}
    \item Each architecture defines a system-call number for each service.
    \item The numbers are part of the \key{application binary interface}.
    \item A kernel cannot renumber existing calls without breaking every
      compiled program.
    \item New calls are added at the end of the table; old calls remain for
      compatibility.
  \end{itemize}
  \vspace{0.2cm}
  \begin{block}{Consequence}
    The system-call table is a long-lived contract between the kernel and every
    binary that runs on it.
  \end{block}
\end{frame}

\begin{frame}{The Cost of Crossing the Boundary}
  \begin{itemize}
    \item Save user context and switch to kernel mode.
    \item Index the table and enter the handler.
    \item Validate arguments and copy data across the boundary.
    \item Perform the operation, then restore the user context.
  \end{itemize}
  \vspace{0.2cm}
  \begin{exampleblock}{Why batching matters}
    Copying 1\,MB in one \texttt{read} costs one mode switch; copying it in
    1-byte chunks costs a million. The bytes are the same; the crossings are
    not.
  \end{exampleblock}
  \muted{\texttt{strace} shows the actual calls a program makes --- useful for
  seeing the boundary in action.}
\end{frame}

\begin{frame}{Common Misconceptions}
  \begin{columns}[T]
    \begin{column}{0.48\textwidth}
      \begin{block}{Wrong}
        \begin{itemize}
          \item ``\texttt{printf} is a system call.''
          \item ``A system call is just a function call.''
          \item ``The kernel can trust the user's pointer.''
          \item ``Every library call traps into the kernel.''
        \end{itemize}
      \end{block}
    \end{column}
    \begin{column}{0.48\textwidth}
      \begin{block}{Right}
        \begin{itemize}
          \item \texttt{printf} uses buffered \texttt{write} calls.
          \item It changes privilege level and uses a kernel entry point.
          \item User pointers and lengths are untrusted and validated.
          \item Many library calls stay entirely in user space.
        \end{itemize}
      \end{block}
    \end{column}
  \end{columns}
\end{frame}

\begin{frame}{Takeaways}
  \begin{itemize}
    \item A \key{system call} is the program's request to the kernel; it is
      invoked through a trap, not an ordinary call.
    \item The \key{system-call number} selects a handler from the kernel's
      dispatch table.
    \item Arguments travel in \key{registers}, a \key{memory block}, or the
      \key{stack}; every pointer and length is validated.
    \item The six categories --- process control, file management, device
      management, information maintenance, communication, protection ---
      organize the interface.
    \item Mode switches are expensive; batch work to reduce them.
  \end{itemize}
\end{frame}

\begin{frame}{Next: System Programs and OS Structures}
  \begin{itemize}
    \item Today: how a program names a kernel service, how the arguments cross
      the boundary, and what kinds of calls exist.
    \item Next lecture: the user-level programs built on top of these calls,
      and the internal structures --- monolithic, layered, microkernel,
      modular, virtual machine --- that organize the kernel itself.
  \end{itemize}
  \vspace{0.3cm}
  \begin{block}{Reading}
    Silberschatz et al., \emph{Operating System Concepts}, Ch.~2
    \S2.5--2.8 \cite{SilberschatzOSC_operating_system_concepts}; Stallings,
    \emph{Operating Systems: Internals and Design Principles}, Ch.~2 and
    Ch.~4 \S4.3 \cite{StallingsOS_internals_design_principles}.
  \end{block}
\end{frame}

\begin{frame}{Further Reading (Page-Cited)}
  \footnotesize
  \begin{itemize}
    \item A.~Silberschatz, P.~B.~Galvin, G.~Gagne, \emph{Operating System
      Concepts}, 8e --- Ch.~2 \S2.3 (system calls, pp.~55--58), \S2.4 (types
      of system calls, pp.~58--66).
    \item D.~P.~Bovet \& M.~Cesati, \emph{Understanding the Linux Kernel}, 3e
      --- Ch.~10 (system calls, pp.~398--419), for the concrete x86/Linux
      entry and exit path.
    \item W.~Stallings, \emph{Operating Systems: Internals and Design
      Principles}, 6e --- Ch.~3 (process description and control), for the
      process context in which calls execute.
  \end{itemize}
  \vspace{0.2cm}
  \muted{Page-verified via each book's \texttt{index.md} under
  \texttt{master\_supporting\_docs/CS502/supporting\_books/}.}
\end{frame}

\begin{frame}[allowframebreaks]{References}
  \footnotesize
  \bibliographystyle{plain}
  \bibliography{../../Bibliography_base}
\end{frame}

\end{document}
